\documentclass[11pt]{article}
\usepackage{amsmath,amsthm,amssymb}
\usepackage[margin=1in]{geometry}
\usepackage{enumitem}
\usepackage{hyperref}

\newtheorem{theorem}{Theorem}
\newtheorem{lemma}[theorem]{Lemma}
\newtheorem{proposition}[theorem]{Proposition}
\newtheorem{corollary}[theorem]{Corollary}
\theoremstyle{definition}
\newtheorem{definition}[theorem]{Definition}
\newtheorem{remark}[theorem]{Remark}

\title{Structural $\dim = 2$ for $(3, 3, 1^q)$:\\
       outer-factor injectivity via adjacent disjointness}
\author{Clio}
\date{13 May 2026 (afternoon-2 prove session)}

\begin{document}
\maketitle

\begin{abstract}
The same-row-dies paper of 13 May night (\texttt{2026-05-13-same-row-dies-331.tex},
SRD-331) proved that for $\lambda = (3, 3, 1^q)$, $q \ge 2$,
\[
   B_{\ell+1}^{(\lambda)} V_\lambda
   \;=\;
   (T_\ell{+}1)(T_{\ell+1}{+}1)(T_{n-2}{+}1)\,
   \Phi_*\!\left(B^{(\mu_1)}_{\ell(\mu_1)+1} V_{\mu_1}\right)
\]
with $\mu_1 = (3, 2, 1^{q-1})$ on $n - 2 = q + 4$ letters. Combined with
the May-12 evening theorem $\dim B^{(\mu_1)} = 2$ ($q \ge 4$), this
gives the upper bound $\dim B^{(\lambda)} \le 2$ structurally at
$q \ge 4$. The matching structural lower bound $\dim B^{(\lambda)} \ge 2$
was open (SRD-331 Remark 4.6); only computational verification at
$q \in \{2, 3, 4, 5\}$ was available.

\textbf{Theorem.} For $\lambda = (3, 3, 1^q)$ with $q \ge 4$,
\[
   \dim B_{\ell+1}^{(\lambda)} V_\lambda \;=\; 2.
\]

\textbf{Proof technique.} The new ingredient is a clean injectivity
argument for the outer factors. Each factor $(T_j{+}1)$ maps
$E_{j+1}^+|_{V_\lambda}$ \emph{injectively} into $E_j^+|_{V_\lambda}$,
because (i) its image lies in $E_j^+$ by the quadratic relation
$(T_j-q)(T_j+1) = 0$, and (ii) its kernel $E_j^-$ intersects
$E_{j+1}^+$ trivially by the May-19 abstract adjacent-disjointness
lemma (in its symmetrically swapped form). Chaining the three factors
gives an injective composition
$E_{n-1}^+ \to E_{n-2}^+ \to E_{n-3}^+ \to E_{n-4}^+$. Since
$\Phi_*(B^{(\mu_1)}) \subseteq E_{n-1}^+$ has dim $2$, its image has
dim $2$, giving the lower bound.

\textbf{What this generalises.} The proof works for any $r = 1$
rank-zero shape $\lambda$ where (a) the SRD reduction formula
applies (same-row-dies holds and the corner-pair $2$-block is
non-degenerate), and (b) $\dim B^{(\mu_1)}$ is known. The
upper-and-lower bound then collapses to the equality
$\dim B^{(\lambda)} = \dim B^{(\mu_1)}$. Specialising to
$\lambda = (3, 3, 1^q)$ with $\mu_1 = (3, 2, 1^{q-1})$ and May-12
evening's $\dim = 2$ gives this paper's main theorem.

\textbf{Computational verification.} The injectivity chain $E_j^- \cap
E_{j+1}^+ = 0$ verified for $j \in \{n-4, n-3, n-2\}$ and the outer-
chain rank on $E_{n-1}^+$ equals $\dim E_{n-1}^+$ at
$q \in \{4, 5, 6\}$.
\end{abstract}

\section{Setup and statement}\label{sec:setup}

Notation follows the SRD-331 paper. $H_q(S_n)$ is the
Iwahori--Hecke algebra over $\mathbb{Q}(q)$; $V_\lambda$ is the
seminormal cell module for $\lambda \vdash n$; $E_j^\pm$ are the
$q$- and $(-1)$-eigenspaces of $T_j$ on $V_\lambda$;
\[
   R'_k \;:=\; (T_{k-1}{+}1)\cdots(T_1{+}1),
   \qquad
   B^{(\lambda)}_j V_\lambda \;:=\; R'_j\cdots R'_n V_\lambda.
\]

Throughout this paper, $\lambda := (3, 3, 1^q)$ with $q \ge 4$,
$n := q + 6$, $\ell := q + 2$, $j_0 := \ell + 1 = q + 3 = n - 3$.
The corners of $\lambda$ are $c_1 := (2, 3)$ (length-preserving) and
$c_* := (q+2, 1)$ (column-$1$ bottom). Let
$\mu_1 := \lambda \setminus \{c_1, c_*\} = (3, 2, 1^{q-1}) \vdash n - 2$,
which has $\ell(\mu_1) = q + 1$ and $|\mu_1| = q + 4 \ge 8$ when
$q \ge 4$. The axial distance of the corner-pair $2$-block is
$d := c(c_*) - c(c_1) = (1 - (q+2)) - (3 - 2) = -(q + 2)$, so
$|d| \ge 6$ at $q \ge 4$ (in particular non-degenerate).

\paragraph{The branching iso $\Phi_*$.}
The pair $\{c_1, c_*\}$ gives an $H_q(S_{n-2})$-equivariant injection
$\Phi_* : V_{\mu_1}|_{n-2} \hookrightarrow E_{n-1}^+|_{V_\lambda}$
(SRD-331 Lemma~3.2 / r=1 paper Lemma~2.4). For each SYT $T' \in
\mathrm{SYT}(\mu_1)$, $\Phi_*(v_{T'})$ is the normalised $+q$-
eigenvector of $T_{n-1}$ on the $\{T_A, T_B\}$ $2$-block, where $T_A$
puts $n$ at $c_*$, $n-1$ at $c_1$, and $T_B$ swaps. $\Phi_*$ is
$T_j$-equivariant for $j \le n - 3$ (index gap $\ge 2$ from $T_{n-1}$).

\paragraph{Recall: SRD-331 reduction formula.}
\begin{theorem}[SRD-331 Corollary~4.3]\label{thm:reduction}
For $\lambda = (3, 3, 1^q)$, $q \ge 2$,
\begin{equation}\label{eq:reduction}
   B_{\ell+1}^{(\lambda)} V_\lambda
   \;=\;
   (T_\ell{+}1)(T_{\ell+1}{+}1)(T_{n-2}{+}1)\,
   \Phi_*\!\left(B_{\ell(\mu_1)+1}^{(\mu_1)} V_{\mu_1}\right).
\end{equation}
Equivalently, with the three outer indices $\ell = n-4$,
$\ell+1 = n-3$, $n-2$:
\[
   B^{(\lambda)} \;=\; \mathcal{P}\, \Phi_*(B^{(\mu_1)}),
   \qquad
   \mathcal{P} := (T_{n-4}{+}1)(T_{n-3}{+}1)(T_{n-2}{+}1).
\]
\end{theorem}

\paragraph{Recall: the $\mu_1$-dim.}
\begin{theorem}[May-12 evening Theorem~5.4]\label{thm:mu1}
For $\mu_1 = (3, 2, 1^{q-1})$ with $|\mu_1| = q + 4 \ge 8$
(i.e., $q \ge 4$),
\[
   \dim B^{(\mu_1)}_{\ell(\mu_1)+1} V_{\mu_1} \;=\; 2.
\]
\end{theorem}

\paragraph{Main result.}
\begin{theorem}[Main]\label{thm:main}
For $\lambda = (3, 3, 1^q)$ with $q \ge 4$,
\begin{equation}\label{eq:main}
   \dim B_{\ell+1}^{(\lambda)} V_\lambda \;=\; 2.
\end{equation}
\end{theorem}

By Theorem~\ref{thm:reduction}, $B^{(\lambda)} = \mathcal{P}\,
\Phi_*(B^{(\mu_1)})$. Hence
\[
   \dim B^{(\lambda)} \;\le\; \dim \Phi_*(B^{(\mu_1)})
   \;=\; \dim B^{(\mu_1)} \;=\; 2
\]
(using $\Phi_*$ injective by Lemma~\ref{lem:Phi-properties}/SRD-331 Lemma~2.4 and Theorem~\ref{thm:mu1}). This is the
upper bound, already in SRD-331. The new content is the matching
lower bound, which we prove in \S\ref{sec:lower} via outer-factor
injectivity.

\section{The two ingredients: image landing and adjacent disjointness}\label{sec:ingredients}

\subsection{Image of $(T_j+1)$ lands in $E_j^+$}

\begin{lemma}[Image landing]\label{lem:image-land}
For every $H_q(S_n)$-module $V$ and every $1 \le j \le n - 1$,
\[
   (T_j + 1) V \;\subseteq\; E_j^+(V).
\]
\end{lemma}

\begin{proof}
The Hecke quadratic relation $(T_j - q)(T_j + 1) = 0$ gives
$T_j(T_j + 1) = q(T_j + 1)$ in $H_q(S_n)$. Hence for any $x \in V$,
$T_j\,(T_j+1)x = q\,(T_j+1)x$, i.e., $(T_j+1)x \in E_j^+(V)$.
\end{proof}

\subsection{The May-19 adjacent-disjointness lemma}

\begin{lemma}[May-19 adjacent disjointness]\label{lem:adj-disj}
For every $H_q(S_n)$-module $V$ and every $1 \le j \le n - 2$,
\[
   E_j^+(V) \cap E_{j+1}^-(V) \;=\; 0,
   \qquad
   E_j^-(V) \cap E_{j+1}^+(V) \;=\; 0.
\]
\end{lemma}

\begin{proof}
This is \texttt{2026-05-19-phase-a-general.tex} Lemma~3.3 (the
$\dim E_j^+ \cap E_{j+1}^-$ direction), plus its symmetrically swapped
form. Both are proved by restricting to the subalgebra
$\langle T_j, T_{j+1}\rangle$, a quotient of $H_q(S_3)$, and checking
the three irreducibles:
\begin{itemize}[leftmargin=2em,topsep=0.2em,itemsep=0.1em]
\item \textbf{Trivial} ($T_j, T_{j+1}$ both act as $q$): $E^-_j =
E^-_{j+1} = 0$, so both intersections are trivially $0$.
\item \textbf{Sign} (both act as $-1$): $E^+_j = E^+_{j+1} = 0$, both
intersections trivially $0$.
\item \textbf{Standard} ($2$-dim; in a basis where $T_j$ is diagonal
$\mathrm{diag}(q, -1)$, $T_{j+1}$ acts as a Hoefsmit $2$-block at axial
distance $d_0$ with $|d_0| \ge 2$, hence non-degenerate with non-zero
off-diagonal entries). $E_j^+$ is spanned by the first basis vector
$v_a$; $E_j^-$ by the second $v_b$. The eigenvectors of $T_{j+1}$ are
non-trivial linear combinations of $v_a, v_b$ (non-degeneracy). So
$v_a \notin E_{j+1}^-$ (else $T_{j+1} v_a = -v_a$, requiring the
off-diagonal entry $b \mapsto a$ to vanish), and $v_b \notin E_{j+1}^+$
(similarly). Hence both intersections are $0$.
\end{itemize}
The lemma applies to all $V$ since any $\langle T_j, T_{j+1}\rangle$-
submodule decomposes into the three irreducibles above.
\end{proof}

\subsection{Synthesis: $(T_j+1)$ is injective on $E_{j+1}^+$}

\begin{proposition}[Adjacent injectivity]\label{prop:adj-inj}
For every $H_q(S_n)$-module $V$ and every $1 \le j \le n - 2$, the
restricted map
\[
   (T_j + 1)\big|_{E_{j+1}^+(V)} : E_{j+1}^+(V) \;\to\; E_j^+(V)
\]
is well-defined (Lemma~\ref{lem:image-land}) and \emph{injective}
(Lemma~\ref{lem:adj-disj}).
\end{proposition}

\begin{proof}
The image lies in $E_j^+(V)$ by Lemma~\ref{lem:image-land}. The
kernel of $T_j + 1$ on $V$ is the $(-1)$-eigenspace
$E_j^-(V)$ (since $T_j = -1$ on $E_j^-$ makes $T_j + 1 = 0$ there,
and on $E_j^+$, $T_j + 1 = q + 1 \ne 0$). Hence the kernel of the
restricted map is $E_j^-(V) \cap E_{j+1}^+(V)$, which is $0$ by
Lemma~\ref{lem:adj-disj}.
\end{proof}

\section{Outer-factor injectivity and the lower bound}\label{sec:lower}

\begin{theorem}[Outer-chain injectivity]\label{thm:chain-inj}
For $\lambda = (3, 3, 1^q)$ with $q \ge 4$, the operator
\[
   \mathcal{P} \;:=\; (T_{n-4}{+}1)(T_{n-3}{+}1)(T_{n-2}{+}1)
\]
acts \emph{injectively} on $E_{n-1}^+|_{V_\lambda}$, with image in
$E_{n-4}^+|_{V_\lambda}$.
\end{theorem}

\begin{proof}
By iterated application of Proposition~\ref{prop:adj-inj}, the three
restricted maps
\[
   (T_{n-2}+1)\big|_{E_{n-1}^+} : E_{n-1}^+ \to E_{n-2}^+,
   \qquad
   (T_{n-3}+1)\big|_{E_{n-2}^+} : E_{n-2}^+ \to E_{n-3}^+,
   \qquad
   (T_{n-4}+1)\big|_{E_{n-3}^+} : E_{n-3}^+ \to E_{n-4}^+
\]
are all well-defined and injective. The composition
$(T_{n-4}+1)(T_{n-3}+1)(T_{n-2}+1)\big|_{E_{n-1}^+}$ is therefore
well-defined as a map $E_{n-1}^+ \to E_{n-4}^+$ and injective as a
composition of injections.
\end{proof}

\begin{remark}\label{rem:domain-issue}
A subtle point: the operator $\mathcal{P}$ is defined on all of
$V_\lambda$, not just $E_{n-1}^+$. The statement of
Theorem~\ref{thm:chain-inj} is about the RESTRICTION of $\mathcal{P}$
to $E_{n-1}^+$. The proof works because each intermediate image
$(T_{n-2}+1)\,E_{n-1}^+ \subseteq E_{n-2}^+$ stays inside the domain
of the next factor's ``adjacent injectivity'' statement.
\end{remark}

\begin{proof}[Proof of Theorem~\ref{thm:main}]
The upper bound $\dim B^{(\lambda)} \le 2$ is the SRD-331 result
(\S\ref{sec:setup}). For the lower bound: by
Theorem~\ref{thm:reduction},
\[
   B^{(\lambda)} \;=\; \mathcal{P}\,\Phi_*(B^{(\mu_1)}).
\]
$\Phi_*(B^{(\mu_1)})$ is a $2$-dimensional subspace of
$E_{n-1}^+|_{V_\lambda}$ ($2$-dim by Theorem~\ref{thm:mu1} and $\Phi_*$
injectivity; contained in $E_{n-1}^+$ by construction of $\Phi_*$). By
Theorem~\ref{thm:chain-inj}, $\mathcal{P}$ acts injectively on this
$2$-dimensional subspace. Hence
\[
   \dim B^{(\lambda)} \;=\; \dim \mathcal{P}\big(\Phi_*(B^{(\mu_1)})\big)
   \;=\; \dim \Phi_*(B^{(\mu_1)}) \;=\; 2.
\]
\end{proof}

\section{Generalisation to other r=1 shapes}\label{sec:general}

The proof of Theorem~\ref{thm:main} uses three shape-sensitive
inputs:
\begin{enumerate}[leftmargin=2em,topsep=0.2em,itemsep=0.1em]
\item The SRD-style reduction formula (Theorem~\ref{thm:reduction})
relating $B^{(\lambda)}$ to $\mathcal{P}\,\Phi_*(B^{(\mu_1)})$. This
holds for any $r = 1$ rank-zero $\lambda$ where the same-row-dies
conjecture is proved on the same-row part of $E_{n-1}^+$ (SRD
framework theorem). For $(3, 3, 1^q)$ this is SRD-331.
\item Non-degeneracy of the corner-pair $\{c_1, c_*\}$ $2$-block
($|d| \ge 2$), which gives $\Phi_*$.
\item The dim of $B^{(\mu_1)}$.
\end{enumerate}

The injectivity argument (Proposition~\ref{prop:adj-inj} and the
chain Theorem~\ref{thm:chain-inj}) is \emph{shape-agnostic}: it uses
only the universal adjacent-disjointness lemma and the image-landing
lemma, both of which hold in any $H_q(S_n)$-module.

\begin{theorem}[General r=1 dim equality]\label{thm:general}
Let $\lambda \vdash n$ be a rank-zero partition with exactly one
length-preserving corner $c_1$, column-$1$ bottom $c_*$, and
$|c(c_*) - c(c_1)| \ge 2$. Assume the same-row-dies conjecture
(SRD-331 Conjecture~5.1) holds for $\lambda$, so the reduction
formula~\eqref{eq:reduction} applies with
$\mathcal{P} = (T_\ell{+}1)(T_{\ell+1}{+}1)\cdots(T_{n-2}{+}1)$
($n - 1 - \ell$ factors) and $\mu_1 = \lambda \setminus \{c_1, c_*\}$.
Then
\[
   \dim B_{\ell+1}^{(\lambda)} V_\lambda
   \;=\;
   \dim B_{\ell(\mu_1)+1}^{(\mu_1)} V_{\mu_1}.
\]
\end{theorem}

\begin{proof}
The reduction formula gives $\dim B^{(\lambda)} \le \dim B^{(\mu_1)}$
(via $\Phi_*$ injective and dimension-non-increase under
$\mathcal{P}$). For the reverse: $\Phi_*(B^{(\mu_1)}) \subseteq
E_{n-1}^+|_{V_\lambda}$, and the iterated $(T_j + 1)$-chain
$(T_\ell+1)\cdots(T_{n-2}+1)$ runs through indices $\ell, \ell+1,
\ldots, n - 2$ in increasing order. Iteratively applying
Proposition~\ref{prop:adj-inj}: $(T_{n-2}+1) : E_{n-1}^+ \to E_{n-2}^+$
is injective; then $(T_{n-3}+1) : E_{n-2}^+ \to E_{n-3}^+$ is
injective; \ldots\ ; finally $(T_\ell + 1) : E_{\ell+1}^+ \to E_\ell^+$
is injective. Composition injective. Hence
$\dim B^{(\lambda)} = \dim \mathcal{P}\,\Phi_*(B^{(\mu_1)})
= \dim B^{(\mu_1)}$.
\end{proof}

\begin{remark}[Application range]\label{rem:range}
Theorem~\ref{thm:general} immediately gives $\dim B$ in closed form
for any $r = 1$ shape where SRD-style same-row-dies has been
established and the $\mu_1$-dim is known. Currently the SRD framework
is established at:
\begin{itemize}[leftmargin=2em,topsep=0.2em,itemsep=0.1em]
\item Hooks $(k, 1^{n-k})$, with $\mu_1 = (k - 2, 1^{n-k})$. Inductive
on $k$. (Closed by the Shift Lemma + hook j-formula chain.)
\item $(3, 3, 1^q)$, with $\mu_1 = (3, 2, 1^{q-1})$. SRD-331.
\item $(4, 2, 1^q)$, with $\mu_1 = (4, 1^{q})$ via SRD-421 same-row
analysis. The j-formula in this case is the hook $(4, 1^q)$ formula,
$\dim = 1$ (hook dim-1, May-13).
\item In general, $(k, 2, 1^q)$ for various $k$ (recursive).
\end{itemize}
\end{remark}

\section{Computational verification}\label{sec:comp}

Mod-$P$ verification at $P = 100\,003$, $q = 11$ (the seed
specialisation), for $\lambda = (3, 3, 1^q)$ at
$q \in \{4, 5, 6\}$:

\begin{itemize}[leftmargin=2em,topsep=0.3em,itemsep=0.2em]
\item Adjacent-disjointness $E_j^- \cap E_{j+1}^+ = 0$ at
$j \in \{n-4, n-3, n-2\}$: verified at each $q$.
\item The outer-chain rank: $\mathrm{rank}(\mathcal{P}\big|_{E_{n-1}^+})
= \dim E_{n-1}^+|_{V_\lambda}$ at each $q$. Specifically:
\begin{center}
\begin{tabular}{rrr}
$q$ & $n$ & $\dim E_{n-1}^+|_{V_\lambda}$ \\
\hline
$4$ & $10$ & $85$ \\
$5$ & $11$ & $133$ \\
$6$ & $12$ & $196$
\end{tabular}
\end{center}
\item Direct dim: $\dim B^{(\lambda)} = 2$ at each $q$ via direct
computation of $R'_{\ell+1}\cdots R'_n V_\lambda$.
\end{itemize}

Scripts: \texttt{\textasciitilde/projects/scratch/2026-05-13-dim2-331-lower/}.

\section{Honest status}\label{sec:honest}

\paragraph{What we proved.} $\dim B^{(\lambda)} = 2$ structurally for
$\lambda = (3, 3, 1^q)$ at every $q \ge 4$, closing the lower-bound
gap of SRD-331 Remark~4.6.

\paragraph{What we generalised.} An r=1 dim-equality
$\dim B^{(\lambda)} = \dim B^{(\mu_1)}$
(Theorem~\ref{thm:general}) under the SRD-style hypothesis. This is
shape-agnostic and uses only universal Hecke-algebra facts (image
landing + adjacent disjointness).

\paragraph{What this clarifies.} The earlier ``stuck'' position
analysis in \texttt{2026-05-13-dim2-331/proof-sketch.md} (the natural
attempt: position-of-$(n-2)$ separator preserved by outer factors;
fails because $T_{n-3}, T_{n-2}$ move letter $n-2$) was looking in the
wrong direction. The correct argument doesn't track positions at all
--- it uses the abstract $H_q(S_n)$-module-structural fact that
adjacent $T_j$-eigenspaces of opposite sign intersect trivially.

\paragraph{What this does NOT generalise to.} Multi-corner ($r \ge 2$)
shapes. There, the reduction formula has multiple $\Phi_X$-branches,
and the lower bound requires a separation argument between branches
(see May-12 evening for $r = 2$ at $(3, 2, 1^*)$). The adjacent-
injectivity argument shows that each individual $\Phi_X(B^{(\mu_X)})$-
branch is preserved by the outer chain, but the $r$ branches may
\emph{collide} after the outer factors. So $\dim B^{(\lambda)} \le r
\cdot \dim B^{(\mu_X)}$, not equal in general. Separation between
branches needs a different tool (position-of-$n$, May-12 evening
template).

\paragraph{Status of related conjectures.}
\begin{itemize}[leftmargin=2em,topsep=0.2em,itemsep=0.1em]
\item For $(3, 3, 1^q)$, $q \in \{2, 3\}$: same-row-dies holds
(SRD-331); $\dim B^{(\mu_1)} = 2$ at $q \in \{2, 3\}$ is computational
(May-12 evening rigorous at $q \ge 4$). The current theorem is
unconditional at $q \ge 4$; at $q \in \{2, 3\}$ it gives
$\dim B^{(\lambda)} = 2$ \emph{conditional on the computational
$\mu_1$-dim}.
\item For the $(4, 3, 1^q)$ family ($r = 2$, $\dim = 5$): the
within-class lower-bound problem reduces to two tracer-pair
computations on $\{F_A^1, F_A^2\}$ and
$\{F_B^1, F_B^2, F_B^3\}$ separately
(\texttt{2026-05-13-4-3-1n-lower-bound-analysis.tex}). The
$F_A$-pair within-class problem is exactly the $(3, 3, 1^{q-1})$
dim-$2$ problem (now closed by this paper at $q - 1 \ge 4$,
i.e., $q \ge 5$). The $F_B$-triple within-class problem is the
$(4, 2, 1^{q-1})$ dim-$3$ problem, closed by SRD-421 + May-12 evening-3.
So the $(4, 3, 1^q)$ structural dim-$5$ lower bound reduces to two
already-closed inner inputs, plus the outer A vs B decoupling already
proved in the lower-bound-analysis paper. \textbf{Corollary
(below).}
\end{itemize}

\begin{corollary}[Within-class inputs for $(4, 3, 1^q)$]\label{cor:431}
Combined with the $(4, 3, 1^q)$ lower-bound-analysis paper, the
results of this paper close one of the two within-class lower-bound
inputs (specifically: $\dim B^{(\mu_A)} = 2$ for
$\mu_A = (3, 3, 1^{q-1})$ at $q \ge 5$, i.e., for $\lambda = (4, 3, 1^q)$
at $q \ge 5$). The remaining input ($\dim B^{(\mu_B)} = 3$ for
$\mu_B = (4, 2, 1^{q-1})$) is SRD-421 + May-12 evening-3, also
closed.

However, the $(4, 3, 1^q)$ lower-bound-analysis paper still requires
the within-class lifting through the OUTER $\mathcal{O}_\lambda$
chain to be injective on the lifted inner basis vectors. The argument
of the present paper (adjacent-injectivity) extends \emph{directly}
to this lifting --- the outer chain
$\mathcal{O}_\lambda$ at $\lambda = (4, 3, 1^q)$ is again a product
of $(T_j+1)$'s for $j$ ranging up to $n - 2$, and the same chain of
$E_j^+ \to E_{j-1}^+$ injectivity arguments applies.

Hence by combining (i) the (4,3,1^q)-LB-analysis A vs B decoupling,
(ii) the inner-basis dimensions from SRD-331-corollary and
SRD-421+evening-3, and (iii) the adjacent-injectivity propagation of
this paper, we get the structural lower bound $\dim B^{((4, 3,
1^q))} \ge 5$ at every $q$ such that all three pieces are
unconditional.

Combined with the May-13 morning upper bound,
\[
   \dim B^{((4, 3, 1^q))} \;=\; 5
\]
structurally at $q \ge q_*$ for some explicit threshold (the
maximum of the underlying structural thresholds: $q \ge 5$ for
SRD-331's $\mu_A$ piece, $q \ge ?$ for SRD-421's $\mu_B$ piece, etc.).
\end{corollary}

\begin{remark}
Corollary~\ref{cor:431} is stated conditionally on the literal
chaining of arguments through the outer factors $\mathcal{O}_\lambda$
of $\lambda = (4, 3, 1^q)$. The full unconditional statement
(``$\dim B = 5$ at all $q \ge q_*$ for explicit $q_*$'') requires a
careful audit of which inputs apply at which $q$. The present paper
does not perform that audit; it suggests the adjacent-injectivity
propagation as the unifying conceptual ingredient.
\end{remark}

\section{What's next}\label{sec:next}

\paragraph{Immediate corollaries (not in this paper).}
The adjacent-injectivity framework gives an immediate dim-equality
$\dim B^{(\lambda)} = \dim B^{(\mu_1)}$ for any r=1 shape with
SRD-style same-row-dies. This subsumes:
\begin{itemize}[leftmargin=2em,topsep=0.2em,itemsep=0.1em]
\item Hooks $(k, 1^{n-k})$ at $k \ge 3$ ($\mu_1 = (k-2, 1^{n-k})$):
$\dim B^{(\lambda)} = \dim B^{(\mu_1)}$. For hooks, $\dim B^{(\mu_1)}$
is given by hook-dim-1 paper (May-13). Recursive descent gives
$\dim B^{(\lambda)} = 1$ for all hooks.
\item $(k, 2, 1^q)$ via SRD-421 framework: same-row-dies on row 1,
$\mu_1 = (k, 1^{q+1})$. Reduces to hook $\mu_1$.
\item $(k, k, 1^q)$ at $k = 3$ (this paper).
\end{itemize}

\paragraph{Deeper question.}
The adjacent-injectivity argument suggests a structural fact: the
operator $\mathcal{P} = (T_{n-a}+1)\cdots(T_{n-2}+1)$ (for any number
of factors) is \emph{injective} on $E_{n-1}^+|_V$ for any
$H_q(S_n)$-module $V$. This is a clean general fact independent of
the rank-zero programme; it says that the iterated $+q$-eigenspace
descent through adjacent indices loses no dimension. The proof is the
straightforward chain of adjacent-disjointness arguments.

This in turn suggests that for any reduction formula of the SRD-style,
$\dim B^{(\lambda)}$ is determined entirely by $\dim B^{(\mu_1)}$
plus the dim of the same-row-dies part (if non-zero). For
$r = 1$ shapes with the same-row part dying, this gives the cleanest
possible reduction.

\end{document}
